% Adición junto a eq:alpha-recurrencia o inc-ampl-normalizacion.
% No cambia ni elimina ninguna fórmula ni renombra sus símbolos.
\paragraph{Bilan signé et somme avec quotient entrant.}
\label{ampl-alfa-z-s-aclaracion}
Il convient de distinguer les deux étapes consécutives de la division
positionnelle au moyen d'une notation commune. Le bilan antérieur au transport
du quotient est
\[
 Z_j:=P_j+E_j-\Phi_j-K_j^{\mathrm{per}}.
\]
La variable $s_j$ de \eqref{eq:alpha-recurrencia} désigne la somme qui
inclut déjà la contribution entrante, de sorte que
\[
 s_j=Z_j+c_{j+1},\qquad
 A_j=Z_j+c_{j+1}-1000c_j.
\]
Dans \eqref{inc-ampl-normalizacion}, la lettre $s_j$ désigne, localement,
le bilan antérieur à cette contribution, c'est-à-dire $Z_j$ dans la notation
présente. Les deux formules expriment donc la même division
euclidienne avec des étapes nominales différentes. La configuration incidencielle
$H^-_\alpha$ utilise les signes du bilan $Z_j$ de la fenêtre initiale ;
la détermination des chiffres utilise en outre les quotients de liaison
$c_{j+1}$ et $c_j$. Cette distinction conserve conjointement le support
signé et la normalisation du développement.
